% Bagian: first-order-logic
% Bab: syntax-and-semantics
% Subbagian: formation-sequences

\documentclass[../../../include/open-logic-section]{subfiles}

\begin{document}

\olfileid[id]{pl}{syn}{fseq}

\olsection{Barisan Pembentukan}

Mendefinisikan !!{formula}s melalui suatu definisi induktif, beserta teknik
pelengkap berupa pembuktian sifat-sifat !!{formula}s melalui induksi,
merupakan pendekatan yang anggun dan efisien. Namun, pendekatan dari bawah
ke atas dan langkah demi langkah terhadap pembangunan !!{formula}s juga
dapat berguna. Di sini kita melakukannya dengan menggunakan gagasan
\emph{barisan pembentukan}.

\begin{defn}[Barisan pembentukan untuk formula]
\ollabel{defn:fseq-frm}
Suatu barisan berhingga $\tuple{!A_0,\dotsc,!A_n}$ dari untai simbol dalam
bahasa~$\Lang L_0$ merupakan \emph{barisan pembentukan} untuk $!A$ jika
$!A \ident !A_n$ dan, untuk setiap $i \leq n$, $!A_i$ merupakan formula
atomik atau terdapat $j,k < i$ sedemikian sehingga salah satu dari hal
berikut berlaku:
\begin{enumerate}
    \tagitem{prvNot}{$!A_i \ident \lnot !A_j$.}{}%
    \tagitem{prvAnd}{$!A_i \ident (!A_j \land !A_k)$.}{}%
    \tagitem{prvOr}{$!A_i \ident (!A_j \lor !A_k)$.}{}%
    \tagitem{prvIf}{$!A_i \ident (!A_j \lif !A_k)$.}{}%
    \tagitem{prvIff}{$!A_i \ident (!A_j \liff !A_k)$.}{}%
\end{enumerate}
\end{defn}

\begin{ex}
\[
\tuple{
    \Obj p_0,
    \Obj p_1,
    (\Obj p_1 \land \Obj p_0),
    \lnot (\Obj p_1 \land \Obj p_0)
}
\]
merupakan barisan pembentukan untuk
$\lnot (\Obj p_1 \land \Obj p_0)$, demikian pula
\[
\tuple{
    \Obj p_0,
    \Obj p_1,
    \Obj p_0,
    (\Obj p_1 \land \Obj p_0),
    (\Obj p_0 \lif \Obj p_1),
    \lnot (\Obj p_1 \land \Obj p_0)
}.
\]
%
Sebagaimana tampak dari contoh kedua, barisan pembentukan dapat memuat
`sampah': formula yang mubazir atau tidak turut membangun hasil akhirnya.
\end{ex}

\begin{prop}
\ollabel{prop:formed}
Setiap !!{formula}~$!A$ dalam~$\Frm[L_0]$ memiliki barisan pembentukan.
\end{prop}

\begin{proof}
Andaikan $!A$ atomik. Maka barisan~$\tuple{!A}$ merupakan barisan
pembentukan untuk~$!A$.
%
Sekarang andaikan $!B$ dan~$!C$ masing-masing memiliki barisan pembentukan
$\tuple{!B_0,\dotsc,!B_n}$ dan $\tuple{!C_0,\dotsc,!C_m}$.
%
\begin{enumerate}
  \tagitem{prvNot}{Jika $!A \ident \lnot !B$,
      maka $\tuple{!B_0,\dotsc,!B_n,\lnot !B_n}$
      merupakan barisan pembentukan untuk~$!A$.}{}
  \tagitem{prvAnd}{Jika $!A \ident (!B \land !C)$,
      maka $\tuple{!B_0,\dotsc,!B_n,!C_0,\dotsc,!C_m,(!B_n \land !C_m)}$
      merupakan barisan pembentukan untuk~$!A$.}{}
  \tagitem{prvOr}{Jika $!A \ident (!B \lor !C)$,
      maka $\tuple{!B_0,\dotsc,!B_n,!C_0,\dotsc,!C_m,(!B_n \lor !C_m)}$
      merupakan barisan pembentukan untuk~$!A$.}{}
  \tagitem{prvIf}{Jika $!A \ident (!B \lif !C)$,
      maka $\tuple{!B_0,\dotsc,!B_n,!C_0,\dotsc,!C_m,(!B_n \lif !C_m)}$
      merupakan barisan pembentukan untuk~$!A$.}{}
  \tagitem{prvIff}{Jika $!A \ident (!B \liff !C)$,
      maka $\tuple{!B_0,\dotsc,!B_n,!C_0,\dotsc,!C_m,(!B_n \liff !C_m)}$
      merupakan barisan pembentukan untuk~$!A$.}{}
\end{enumerate}
Menurut prinsip induksi pada !!{formula}s, setiap !!{formula} memiliki barisan
pembentukan.
\end{proof}

Kita juga dapat membuktikan arah sebaliknya. Hal ini penting karena
menunjukkan bahwa kedua cara kita mendefinisikan formula adalah ekuivalen:
keduanya memberikan hasil yang sama. Hal ini juga berarti bahwa kita dapat
membuktikan teorema tentang formula dengan menggunakan induksi biasa pada
panjang barisan pembentukan.

\begin{lem}
\ollabel{lem:fseq-init}
Andaikan $\tuple{!A_0,\dotsc,!A_n}$ merupakan barisan pembentukan
untuk~$!A_n$, dan $k \leq n$. Maka $\tuple{!A_0,\dotsc,!A_k}$ merupakan
barisan pembentukan untuk~$!A_k$.
\end{lem}

\begin{proof}
Latihan.
\end{proof}

\begin{thm}
\ollabel{thm:fseq-frm-equiv}
$\Frm[L_0]$ merupakan himpunan semua untai simbol dalam bahasa~$\Lang L_0$
yang memiliki barisan pembentukan.
\end{thm}

\begin{proof}
Misalkan $F$ ialah himpunan semua untai simbol dalam bahasa~$\Lang L_0$ yang
memiliki barisan pembentukan. Kita telah melihat dalam
\olref[pl][syn][fseq]{prop:formed} bahwa $\Frm[L_0] \subseteq F$, sehingga
sekarang kita membuktikan arah sebaliknya.

Andaikan $!A$ memiliki barisan pembentukan $\tuple{!A_0,\dotsc,!A_n}$. Kita
membuktikan bahwa $!A \in \Frm[L_0]$ melalui induksi kuat pada~$n$, yaitu
indeks terakhir barisan tersebut. Hipotesis induksi kita menyatakan bahwa
setiap untai simbol dengan barisan pembentukan berindeks terakhir $m < n$
berada dalam~$\Frm[L_0]$. Menurut definisi barisan pembentukan, $!A_n$
merupakan formula atomik atau harus terdapat $j,k < n$ sedemikian sehingga
salah satu dari hal berikut berlaku:
\begin{enumerate}
    \tagitem{prvNot}{$!A_n \ident \lnot !A_j$.}{}%
    \tagitem{prvAnd}{$!A_n \ident (!A_j \land !A_k)$.}{}%
    \tagitem{prvOr}{$!A_n \ident (!A_j \lor !A_k)$.}{}%
    \tagitem{prvIf}{$!A_n \ident (!A_j \lif !A_k)$.}{}%
    \tagitem{prvIff}{$!A_n \ident (!A_j \liff !A_k)$.}{}%
\end{enumerate}
Sekarang kita bernalar menurut kasus. Jika $!A_n$ atomik, maka
$!A_n \in \Frm[L_0]$. Sebaliknya, andaikan $!A \ident
(!A_j \land !A_k)$. Menurut \olref[pl][syn][fseq]{lem:fseq-init},
$\tuple{!A_0,\dotsc,!A_j}$ dan $\tuple{!A_0,\dotsc,!A_k}$ masing-masing
merupakan barisan pembentukan untuk $!A_j$ dan~$!A_k$. Karena keduanya
merupakan subbarisan awal sejati dari barisan pembentukan untuk~$!A$, indeks
terakhir keduanya lebih kecil daripada~$n$. Oleh karena itu, menurut
hipotesis induksi, $!A_j$ dan~$!A_k$ berada dalam~$\Frm[L_0]$, sehingga
menurut definisi !!a{formula}, $(!A_j \land !A_k)$ juga berada di dalamnya.
Kasus-kasus lain mengikuti penalaran yang sejajar.
\end{proof}

\end{document}
